\textbf{3-1}\quad 在 $p$ 表象求解 $\delta$ 势阱 $V(x) = -\gamma \delta (x)$ 中的定态能量和波函数 $(\gamma > 0)$ .

\vfill

\newpage
\vfill

\noindent
\textbf{3-2}\quad 已知在 $\delta$ 势阱 $V(x) = -\gamma \delta (x)$ 中的定态归一化波函数（ $p$ 表象）为 $\varphi(p) = \frac{A}{p^2 + \hbar^2 k^2}, A = \sqrt{\frac{2\hbar^3 k^3}{\pi}}, k = \frac{\mu\gamma}{\hbar^2}$ ，试计算 $\Delta x \Delta p$ ，验证测不准关系。

\vfill

\newpage
\vfill

\noindent
\textbf{3-3}\quad 在 p 表象计算一维谐振子定态能量和波函数.

\vfill

\newpage
\vfill

\noindent
\textbf{3-4}\quad 质量为 $\mu$ 的粒子在均匀力场 $f(x) = -F(F > 0)$ 中运动, 运动的范围限制在 $x > 0$ . 给出动量表象中的定态方程, 并求出定态波函数 $\varphi(p)$ .

\vfill

\newpage
\vfill

\noindent
\textbf{3-5}\quad 质量为 $\mu$ 的粒子在均匀力场 $f(x) = -F(F > 0)$ 中运动, $\rho(p, t)$ 为其在动量空间的概率密度, 求 $\partial \rho / \partial t$ 与 $\partial \rho / \partial p$ 的关系.

\vfill

\newpage
\vfill

\noindent
\textbf{3-6}\quad 已知 $t = 0$ 时一维自由粒子波函数在坐标表象和动量表象分别是
$$
\psi (x) = N x \exp \left(- a x ^ {2} + \mathrm{i} \frac {p _ {0} x}{\hbar}\right)
$$
$$
\varphi (p) = C \left(p - p _ {0}\right) \exp \left[ - b \left(p - p _ {0}\right) ^ {2} \right]
$$
其中 $a, b$ 和 $p_0$ 都是已知实数， $N$ 与 $C$ 是归一化常数。试求 $t = 0$ 时和 $t > 0$ 时粒子坐标和动量的平均值 $\overline{x}$ 和 $\overline{p}$ 。

\vfill

\newpage
\vfill

\noindent
\textbf{3-7}\quad 中子 $n$ 和反中子 $\overline{n}$ 的质量都是 $m$ , 它们的态 $|n\rangle$ 和 $|\overline{n}\rangle$ 可以看成是一自由哈密顿量 $\hat{H}_0$ 的简并态: $\hat{H}_0 |n\rangle = mc^2 |n\rangle, \hat{H}_0 |\overline{n}\rangle = mc^2 |\overline{n}\rangle$ . 设有某种相互作用 $\hat{H}'$ 能使中子与反中子相互转变: $\hat{H}' |n\rangle = \alpha |\overline{n}\rangle, \hat{H}' |\overline{n}\rangle = \alpha^* |n\rangle$ , 其中 $\alpha = \alpha^*$ . 试求 $t = 0$ 时刻的一个中子在 $t$ 时刻变成反中子的概率.

\vfill

\newpage
\vfill

\noindent
\textbf{3-8}\quad 有一量子体系, 其态矢空间三维. 选择基矢 $\{|1\rangle, |2\rangle, |3\rangle\}$ . 体系的哈密顿量 $\hat{H}$ 及另外两个力学量 $\hat{A}$ 与 $\hat{B}$ 为
$$
\hat {H} = \hbar \omega_ {0} \left( \begin{array}{c c c} 1 & 0 & 0 \\ 0 & 2 & 0 \\ 0 & 0 & 2 \end{array} \right), A = a \left( \begin{array}{c c c} 1 & 0 & 0 \\ 0 & 0 & 1 \\ 0 & 1 & 0 \end{array} \right), B = b \left( \begin{array}{c c c} 0 & 1 & 0 \\ 1 & 0 & 0 \\ 0 & 0 & 1 \end{array} \right)
$$
设 $t = 0$ 时体系的态矢为 $\left|\psi(0)\right\rangle = \frac{1}{\sqrt{2}}\left|1\right\rangle + \frac{1}{2}\left|2\right\rangle + \frac{1}{2}\left|3\right\rangle$ , 

(1) 在 $t = 0$ 时测量体系能量 $H$ 可得哪些结果？相应概率多大？计算 $H$ 平均值 $\bar{H}$ 及 $\Delta H = \sqrt{H^2 - (\bar{H})^2}$ . 

(2) 如 $t = 0$ 时测量 $A$ ，可能值与相应概率有多大？写出测量后体系的态矢量. 

(3) 计算任意 $t$ 时刻 $A$ 与 $B$ 的平均值 $\overline{A(t)}$ 与 $\overline{B(t)}$ .

\vfill

\newpage
\vfill

\noindent
\textbf{3-9}\quad 厄米算符 $\hat{A}$ 与 $\hat{B}$ 满足 $\hat{A}^2 = \hat{B}^2 = 1$ ，且 $\hat{A}\hat{B} + \hat{B}\hat{A} = 0$ 。求

(1)在 $\hat{A}$ 表象中算符 $\hat{A}$ 与 $\hat{B}$ 的矩阵表示；

(2)在 $\hat{A}$ 表象中算符 $\hat{B}$ 的本征值与本征态矢；

(3)由 $\hat{A}$ 表象到 $\hat{B}$ 表象的幺正变换 $S$ 矩阵，并把 $\hat{B}$ 矩阵对角化。

\vfill

\newpage
\vfill

\noindent
\textbf{3-10}\quad 在 $l = 1$ 的 $(\hat{L}^2, \hat{L}_z)$ 表象中, 基矢为
$$
\varphi_ {1} = Y _ {1 1} (\theta , \varphi), \varphi_ {2} = Y _ {1 0} (\theta , \varphi), \varphi_ {3} = Y _ {1 - 1} (\theta , \varphi)
$$
求 $\hat{L}_x, \hat{L}_y$ 与 $\hat{L}_z$ 的矩阵表示.

\vfill

\newpage
\vfill

\noindent
\textbf{3-11}\quad 已知在 $l = 1$ 的 $\left(\hat{L}^2, \hat{L}_z\right)$ 表象中，
$$
\hat {L} _ {x} = \frac {\hbar}{\sqrt {2}} \left( \begin{array}{l l l} 0 & 1 & 0 \\ 1 & 0 & 1 \\ 0 & 1 & 0 \end{array} \right), \hat {L} _ {y} = \frac {\hbar}{\sqrt {2}} \left( \begin{array}{l l l} 0 & - \mathrm{i} & 0 \\ \mathrm{i} & 0 & - \mathrm{i} \\ 0 & \mathrm{i} & 0 \end{array} \right), \hat {L} _ {z} = \hbar \left( \begin{array}{l l l} 1 & 0 & 0 \\ 0 & 0 & 0 \\ 0 & 0 & - 1 \end{array} \right)
$$

(1) 给出它们的本征值与本征态矢；

(2)给出 $(\hat{L}^2, \hat{L}_z)$ 表象到 $(\hat{L}^2, \hat{L}_x)$ 表象变换的 $S$ 矩阵，并通过 $S$ 矩阵，求出在 $(\hat{L}^2, \hat{L}_x)$ 表象中 $\hat{L}_x, \hat{L}_y$ 与 $\hat{L}_z$ 的矩阵表示，本征值与本征态矢.

\vfill

\newpage
\vfill

\noindent
\textbf{3-12}\quad 有一量子体系处于角动量 $\hat{L}^2$ 与 $\hat{L}_z$ 的共同本征态上, 总角动量平方值为 $2\hbar^2$ . 已知测量 $\hat{L}_y$ 得值为 0 的概率是 $1/2$ , 求测量 $\hat{L}_y$ 得值为 $\hbar$ 的概率.

\vfill

\newpage
\vfill

\noindent
\textbf{3-13}\quad 粒子处于态 $\psi = C\mathrm{e}^{-\alpha r}(x + y + 2z)$ , 其中 $\alpha$ 为正数, $C$ 为归一化常数. 求 $L^2$ 的取值, $L_z$ 的平均值, $L_z = \hbar$ 的概率, $L_x$ 的可能值及相应概率.

\vfill

\newpage
\vfill

\noindent
\textbf{3-14}\quad 体系处于态 $\psi = c_{1}Y_{11} + c_{2}Y_{10}$ ( $|c_1|^2 + |c_2|^2 = 1$ ), 求

(1) $L_{z}$ 的可能值及相应概率; 

(2) $L^{2}$ 的可能值及相应概率; 

(3) $L_{x}$ 的可能值及相应概率.

\vfill

\newpage
\vfill

\noindent
\textbf{3-15}\quad 体系处于态 $\psi = c_{1}Y_{11} + c_{2}Y_{20}$ ( $|c_1|^2 + |c_2|^2 = 1$ ), 求

(1) $L_{z}$ 的可能值及相应概率; 

(2) $L^{2}$ 的可能值及相应概率; 

(3) $L_{x}$ 的可能值及相应概率.


\vfill

\newpage
\vfill

\noindent
\textbf{3-16}\quad 在角动量 $\hat{J}^2$ 与 $\hat{J}_z$ 为对角矩阵的表象, 对 $j = 3 / 2$ , 求 $\hat{J}^2, \hat{J}_x, \hat{J}_y$ 和 $\hat{J}_z$ 的矩阵表示.

\vfill

\newpage
\vfill

\noindent
\textbf{3-17}\quad 在 $j = 3 / 2$ 的 $(\hat{J}^2\hat{J}_z)$ 表象中， $\hat{J}_x$ 的矩阵为
$$
\hat {J} _ {x} = \frac {\hbar}{2} \left( \begin{array}{c c c c} 0 & \sqrt {3} & 0 & 0 \\ \sqrt {3} & 0 & 2 & 0 \\ 0 & 2 & 0 & \sqrt {3} \\ 0 & 0 & \sqrt {3} & 0 \end{array} \right)
$$
其中行和列都是按 $\hat{J}_z$ 的量子数 $m$ 由大到小排列的. 

(1) 求出 $\hat{J}_y$ 的矩阵; 

(2) 求出与 $\hat{J}_y$ 最大本征值相应的本征态, 并说明其中各矩阵元的物理意义.

\vfill

\newpage
\vfill

\noindent
\textbf{3-18}\quad 在由正交基矢 $\{|1\rangle, |2\rangle, |3\rangle\}$ 构成的三维态矢空间中, 哈密顿算符 $\hat{H}$ 与力学量 $\hat{A}$ 的矩阵为
$$
\hat {H} = E _ {0} \left( \begin{array}{c c c} 1 & 0 & 0 \\ 0 & - 1 & 0 \\ 0 & 0 & - 1 \end{array} \right), \hat {A} = a \left( \begin{array}{c c c} 1 & 0 & 0 \\ 0 & 0 & 1 \\ 0 & 1 & 0 \end{array} \right)
$$

(1)证明 $\hat{A}$ 为守恒量；

(2)求出 $\hat{H}$ 与 $\hat{A}$ 的共同本征态矢组.

\vfill

\newpage
\vfill

\noindent
\textbf{3-19}\quad 一个空间转子, 其哈密顿量为 $\hat{H} = \frac{\hat{L}_x^2}{2I_x} + \frac{\hat{L}_y^2}{2I_y} + \frac{\hat{L}_z^2}{2I_z}$ . 转子的轨道角动量量子数 $l = 1$ , $I_x, I_y$ 与 $I_z$ 均为正实数. 

(1) 在角动量表象中求出 $\hat{L}_x, \hat{L}_y$ 与 $\hat{L}_z$ 的矩阵表示; 

(2) 求出 $\hat{H}$ 的本征值.

\vfill

\newpage
\vfill

\noindent
\textbf{3-20}\quad 质量为 $\mu$ 的粒子受到力 $F(\boldsymbol{r}) = -\nabla V(\boldsymbol{r})$ 的作用, 粒子的波函数满足动量空间薛定谔方程
$$
\mathrm{i} \hbar \frac {\partial}{\partial t} \phi (\boldsymbol {p}, t) = \left(\frac {p ^ {2}}{2 \mu} - \alpha \nabla_ {\boldsymbol {p}} ^ {2}\right) \phi (\boldsymbol {p}, t)
$$

其中 $\alpha$ 是实常数. 求 $F(r)$ .

\vfill

\newpage